\subsubsection{DatasetPageLoaded}
\label{sec:DatasetPageLoaded}
\paragraph{Descrizione}
DatasetPageLoaded è il componente che permette di visualizzare le coppie di domande-risposte del dataset selezionato. Presenta le seguenti funzionalità:
\begin{itemize}
    \item Scorrimento della lista di domande-risposte attese;
    \item Paginazione della lista di domande-risposte attese;
    \item Modifica delle domande e risposte attese;
    \item Eliminazione di una coppia domanda-risposta attesa;
    \item Ricerca delle domande e risposte tramite parole chiave;
    \item Aggiunta di una coppia domanda-risposta attesa;
    \item Esecuzione di un test con un LLM;
    \item Visualizzazione delle informazioni di un dataset.
\end{itemize}
\paragraph{SottoComponenti}
DatasetPageLoaded è composto dai seguenti sotto-componenti:
\begin{itemize}
        \item \textbf{ReturnBtn}: pulsante che permette di ritornare nella pagina precedente;
        \item \textbf{DatasetNameText}: testo che mostra il nome del dataset selezionato;
        \item \textbf{SearchInfoBar}: barra di ricerca per le domande e le risposte;
        \item \textbf{SaveDatasetBtn}: pulsante che permette di salvare e/o aggiornare il dataset all'interno della base di dati;
        \item \textbf{QAListView}: lista scorrevole che mostra tutte le coppie domande-risposte impaginate:
        \begin{itemize}
            \item \textbf{QAModificaBtn}: pulsante che permette di modificare la coppia domanda-risposta;
            \item \textbf{QADeleteBtn}: pulsante che permette di eliminare la coppia domanda-risposta.
        \end{itemize}

        \item \textbf{Form}: form che permette all'utente di inserire informazioni riguardanti i campi della domanda e della risposta. 
        \begin{itemize}
            \item \textbf{CreateQABtn}: pulsante che permette di aggiungere la coppia domanda-risposta all'interno del dataset selezionato.
        \end{itemize}

        \item \textbf{TestBtn}: pulsante che permette di eseguire un test con un LLM. 
        \begin{itemize}
            \item \textbf{QAIncompleteLink}: finestra di messaggio di errore che contiene i link a tutte le domande e risposte incomplete e da correggere per poter eseguire il test;
            \item \textbf{ConfigLLMBtn}: finestra di messaggio che contiene una lista degli LLM disponibili da utilizzare per il test.
        \end{itemize}

    \end{itemize}

\paragraph{Tracciamento dei requisiti}
\begin{tabularx}{\textwidth}{|c|c|}
    \hline
    \textbf{ID} & \textbf{Componente} \\
    \hline
        RF0-1.01 & DatasetNameText\\
    \hline
        RF0-1.02 & DatasetNameText\\
    \hline
        RF0-1.03 & QAListView\\
    \hline
        RF0-1.04 & QAListView\\
    \hline
        RF0-1.05 & QAListView\\
    \hline
        RF0-1.06 & SearchInfoBar\\
    \hline
        RF0-1.07 & CreateQABtn\\
    \hline
        RF0-1.08 & TestBtn\\
    \hline
        RF0-1.09 & QAIncompleteLink\\
    \hline
        RF0-1.10 & TestBtn\\
    \hline
        RF0-1.11 & SaveDatasetBtn\\
    \hline
        RF0-1.13 & SaveDatasetBtn\\
    \hline
        RFF-1.01 & ConfigLLMBtn\\ %parte di matteo eghosa
    \hline
        RF0-2.01 & QAListView\\
    \hline
        RF0-2.02 & QAListView\\
    \hline
        RF0-2.03 & QAListView\\
    \hline
        RF0-3.01 & QAListView\\
    \hline
        RF0-3.02 & QAListView\\
    \hline
        RF0-3.03 & QAListView\\
    \hline
        RF0-3.04 & QAModificaBtn\\
    \hline
        RF0-3.05 & QADeleteBtn\\
    \hline
        RF0-3.06 & QADeleteBtn\\
    \hline
        RF0-4.01 & QAListView\\
    \hline
        RF0-4.02 & QAListView\\
    \hline
        RF0-4.03 & QAModificaBtn\\
    \hline
        RF0-4.04 & QAModificaBtn\\
    \hline
        RF0-4.05 & QAModificaBtn\\
    \hline
        RF0-5.01 & QAIncompleteLink\\
    \hline
        RF0-5.02 & QAIncompleteLink\\
    \hline
        RF0-5.03 & QAIncompleteLink\\
    \hline
        RF0-8.08 & ReturnBtn\\
    \hline
\end{tabularx}